\section*{Hoofdstuk \ref{ch:b2:ellingham} --- Ellinghamdiagrammen en metallurgie}

\begin{solution}{exo:b2:ellingham:1}
\ce{4Cu + O2 -> 2Cu2O}; \ce{4/3Al + O2 -> 2/3Al2O3}; \ce{2C + O2 -> 2CO};
\ce{C + O2 -> CO2}.
\end{solution}

\begin{solution}{exo:b2:ellingham:2}
$\Delta_r H^\circ = 2(-350.46) = \qty{-700.92}{kJ}$; $\Delta_r S^\circ = 2(43.65)
- 2(41.63) - 205.15 = \qty{-201.1}{J/K}$. Dus $\Delta_r G^\circ = -700.9 +
0.2011\,T$ (\unit{kJ}, $T$ in \unit{K}), geldig tot het smeltpunt van zink.
\end{solution}

\begin{solution}{exo:b2:ellingham:3}
Bij \qty{1500}{K} ligt de lijn van \ce{Cr2O3} rond \qty{-495}{kJ}. Eronder, dus
in staat om chroomoxide te reduceren: silicium, titaan, aluminium, magnesium,
calcium. Erboven, niet in staat: koper, ijzer, zink. Koolstof, bij \qty{-488}{kJ}, ligt
er net boven: het reduceert chroomoxide pas bij een iets hogere temperatuur,
en levert dan chroom dat beladen is met carbide --- een van de redenen waarom chroom voor
koolstofvrije legeringen aluminothermisch gemaakt wordt.
\end{solution}

\begin{solution}{exo:b2:ellingham:4}
De helling is $-\Delta_r S^\circ$, bepaald door de verandering van de hoeveelheid gas.
\ce{C + O2 -> CO2}: één mol gas geeft er één, $\Delta_r S^\circ \approx 0$,
horizontaal. \ce{2C + O2 -> 2CO}: één geeft er twee, $\Delta_r S^\circ > 0$, de
lijn daalt. \ce{2CO + O2 -> 2CO2}: drie geven er twee, $\Delta_r S^\circ < 0$, de
lijn stijgt, zoals die van de metalen.
\end{solution}

\begin{solution}{exo:b2:ellingham:5}
$\Delta_r G^\circ(\qty{600}{K}) = -181.6 + 600 \times 0.2165 = \qty{-51.7}{kJ}$,
$p_{\ce{O2}} = p^\circ\exp(-51\,710/(8.314 \times 600)) = \qty{3.2e-5}{bar}$. Lucht
bevat \qty{0.21}{bar}, veel meer: kwik wordt bij \qty{600}{K} geoxideerd (langzaam
--- zo werd het oxide voor het eerst bereid), en het oxide ontleedt pas
boven ongeveer \qty{730}{K}.
\end{solution}

\begin{solution}{exo:b2:ellingham:6}
$\Delta_r G^\circ = -1582.3 - (-743.5) = \qty{-838.8}{kJ}$. De reagentia zijn
twee vaste stoffen die alleen aan het oppervlak van hun korrels contact maken, en de reactie
heeft een hoge activeringsenergie: ze moet plaatselijk ontstoken worden (een magnesiumlont);
daarna houdt de warmte die ze vrijgeeft haar op gang.
\end{solution}

\begin{solution}{exo:b2:ellingham:7}
De lijn van \ce{2C + O2 -> 2CO} snijdt die van \ce{2Fe + O2 -> 2FeO} rond
\qty{1040}{K}; daarboven heeft \ce{FeO + C -> Fe + CO} een negatieve
\omterm{def:b2:reaction-free-energy:reaction-gibbs}{standaardreactie-gibbsenergie}.
\end{solution}

\begin{solution}{exo:b2:ellingham:8}
$\Delta_r G^\circ = 2(-209.1) - (-396.0) = \qty{-22.2}{kJ}$, $K^\circ =
\exp(22\,200/(8.314 \times 1100)) = 11.3$. Dan is $x^2/(1 - x) = 11.3$, $x =
0.92$. Bij afkoelen tot \qty{700}{K} schuift het evenwicht terug naar \ce{CO2}
($x = 0.016$): koolstofmonoxide disproportioneert, \ce{2CO -> C + CO2}, en
er zet zich roet af --- langzaam, tenzij een oppervlak zoals ijzer het katalyseert.
\end{solution}

\begin{solution}{exo:b2:ellingham:9}
De lijn van water (drie mol gas geven er twee) stijgt minder steil dan de
lijn \ce{CO}/\ce{CO2}, eveneens drie voor twee maar met een groter entropieverlies. De
twee snijden elkaar rond \qty{1100}{K}: daaronder ligt de lijn \ce{CO}/\ce{CO2} lager en
is koolstofmonoxide de sterkere reductor; daarboven ligt de lijn van water lager en
wint waterstof. Het is dezelfde bewering als de verschuiving van het watergas-evenwicht
\ce{CO + H2O <=> CO2 + H2} naar links bij hoge temperatuur.
\end{solution}

\begin{solution}{exo:b2:ellingham:10}
Per mol \ce{O2}, dus voor \ce{2MgO + Si -> SiO2 + 2Mg(g)}: $\Delta_r
G^\circ = -643.7 - (-845.5) = \qty{+201.8}{kJ}$. Met silicium en siliciumdioxide met
activiteit 1 is $\Delta_r G = \Delta_r G^\circ + RT\ln(p_{\ce{Mg}}/p^\circ)^2$, negatief
wanneer $p_{\ce{Mg}} < p^\circ\exp(-201\,800/(2 \times 8.314 \times 1500)) =
\qty{3e-4}{bar}$. Een vacuümoven pompt de magnesiumdamp weg zo snel
als hij ontstaat en condenseert hem op een koud oppervlak, zodat zijn druk die waarde
nooit bereikt; kalk die aan de lading wordt toegevoegd, bindt het siliciumdioxide en verlaagt de activiteit ervan,
wat nog meer helpt.
\end{solution}

\begin{solution}{exo:b2:ellingham:11}
Vier deeltjes, één reactie, drie fasen (twee vaste stoffen en het gas): $v = 4 - 1 +
2 - 3 = 2$. Bij gegeven temperatuur en druk ligt de verhouding $p_{\ce{CO}}/p_{\ce{CO2}}$
vast: het gas kan niet al zijn koolstofmonoxide aan het ijzeroxide afstaan,
en het topgas bevat er nog veel van --- dat verbrand wordt om de
lucht voor te verwarmen.
\end{solution}

\begin{solution}{exo:b2:ellingham:12}
$\log K = 2(0.34 + 0.41)/0.059 = 25.4$, $K = \num{3e25}$. Met
$[\ce{Fe^2+}] = \qty{1.0}{mol/L}$ is $[\ce{Cu^2+}] = 1/K \approx
\qty{4e-26}{mol/L}$: het koper wordt volledig verwijderd.
\end{solution}

\begin{solution}{pb:b2:ellingham:1}
\textbf{1.} $\Delta_r H^\circ = \qty{-700.92}{kJ}$, $\Delta_r S^\circ =
\qty{-201.1}{J/K}$.
\textbf{2.} $\Delta_r G^\circ = -700.9 + 0.2011\,T$ (\unit{kJ}), onder
\qty{692.7}{K}.
\textbf{3.} $\Delta_{\text{fus}}S = 7320/692.7 = \qty{10.6}{J/(K.mol)}$;
$\Delta_{\text{vap}}S = 115\,300/1180.2 = \qty{97.7}{J/(K.mol)}$.
\textbf{4.} Vloeibaar zink: $\Delta_r H^\circ = -700.92 - 2(7.32) = \qty{-715.6}{kJ}$,
$\Delta_r S^\circ = -201.1 - 2(10.6) = \qty{-222.2}{J/K}$: $\Delta_r G^\circ =
-715.6 + 0.2222\,T$.
\textbf{5.} Gasvormig zink: $\Delta_r H^\circ = -715.6 - 2(115.3) = \qty{-946.2}{kJ}$,
$\Delta_r S^\circ = -222.2 - 2(97.7) = \qty{-417.7}{J/K}$: $\Delta_r G^\circ =
-946.2 + 0.4177\,T$.
\textbf{6.} Hellingen $0.201$, $0.222$ en $0.418$ \unit{kJ/K}. Boven het kookpunt
verbruikt de reactie drie mol gas in plaats van één: het entropieverlies, en dus de helling,
verdubbelt.
\textbf{7.} Bij \qty{692.7}{K}: $-700.9 + 139.3 = -715.6 + 153.9 = \qty{-561.6}{kJ}$;
bij \qty{1180.2}{K}: $-715.6 + 262.3 = -946.2 + 492.9 = \qty{-453.3}{kJ}$. De
lijnen sluiten op elkaar aan, zoals het hoort.
\textbf{8.} Helling $(-453.0 + 418.2)/200 = \qty{-0.174}{kJ/K}$, $a = -418.2 +
0.174 \times 1100 = \qty{-226.8}{kJ}$: $\Delta_r G^\circ = -226.8 - 0.174\,T$.
\textbf{9.} $\Delta_r S^\circ = \qty{+174}{J/K}$: één mol gas geeft er twee.
\textbf{10.} \ce{2ZnO + 2C -> 2Zn(g) + 2CO}, $\Delta_r G^\circ = (-226.8 - 0.174\,T)
- (-946.2 + 0.4177\,T) = 719.4 - 0.592\,T$ (\unit{kJ}).
\textbf{11.} Nul bij $T = 719.4/0.592 = \qty{1215}{K}$.
\textbf{12.} Boven zijn kookpunt ontstaat zink als damp: de reactie maakt
gas uit vaste stoffen, haar entropie is groot en positief, en de damp verlaat de
lading, wat het metaal door destillatie scheidt van het erts en de cokes.
\textbf{13.} Vier deeltjes, één reactie, één bijzondere voorwaarde ($p_{\ce{Zn}} =
p_{\ce{CO}}$), drie fasen: $v = 4 - 1 - 1 + 2 - 3 = 1$. De druk vastleggen
legt de evenwichtstemperatuur vast.
\textbf{14.} Eén mol zinkdamp per mol koolstofmonoxide:
$p_{\ce{Zn}} = p_{\ce{CO}} = \qty{0.5}{bar}$.
\textbf{15.} Per mol \ce{ZnO} is $\Delta_r G^\circ = \frac12(719.4 - 0.592 \times
1300) = \qty{-25.1}{kJ/mol}$, $K^\circ = \exp(25\,100/(8.314 \times 1300)) = 10$;
$Q = 0.25 < K^\circ$: de reductie verloopt.
\textbf{16.} Net boven het snijpunt is de reductie al begunstigd (met
\qty{0.5}{bar} van elk gas vanaf ongeveer \qty{1170}{K}); heter kost brandstof,
verslijt de kleiretorten en levert niet meer zink, want de reactie verloopt
hoe dan ook volledig.
\textbf{17.} \ce{Zn(g) + CO2 -> ZnO + CO}.
\textbf{18.} De lijn van zinkoxide ligt onder de lijn \ce{CO}/\ce{CO2}
($-445$ tegenover \qty{-357}{kJ} bij \qty{1200}{K}): zinkdamp reduceert koolstofdioxide
en wordt weer oxide. In de retort vernietigt de hete koolstof elk spoor
\ce{CO2} (Boudouard); in de koelere condensor doet niets dat, en \ce{CO2} --- uit
binnenlekkende lucht, of uit \ce{2CO -> C + CO2} bij afkoeling --- bederft het zink tot een
grijs geoxideerd poeder.
\textbf{19.} Langzaam afkoelen laat de damp lang in het gebied waar hij
opnieuw oxideert, en geeft fijn stof in plaats van vloeibaar metaal; snelle
condensatie tot vloeistof beperkt beide.
\textbf{20.} $10^6/65.38 = \qty{1.53e4}{mol}$ zink, evenveel koolstof:
$1.53 \times 10^4 \times 12.011 = \qty{184}{kg}$ (in de praktijk veel meer, om de
retorten te verwarmen).
\textbf{21.} $\Delta_r H^\circ = \frac12(-226.8 + 946.2) = \qty{+360}{kJ}$ per
mol zink; per ton $1.53 \times 10^4 \times 360 = \qty{5.5e6}{kJ}$, dus
\qty{5.5}{GJ} --- geleverd door meer steenkool rond de retorten te verbranden.
\textbf{22.} Koolstof reduceert zinkoxide onder standaardomstandigheden boven
$\boldsymbol{T \approx \qty{1.22e3}{K}}$.
\end{solution}
